\scanpage{7}
\mspage{\handwritten{10}}
Lorsque \deleted{\unsure{$X$ et $W$ sont}} de plus \inserted{$W$ est} réduit\deleted{\unsure{s}}, et si $U$ désigne le plus grand ouvert de $V$ tel que $\deleted{\phi}\handwritten{\psi}$ induise un isomorphisme $\deleted{\phi}\handwritten{\psi}^{-1}(U)\handwritten{\to}U$, alors les conditions précédentes équivalent aussi à

(iv) $a\handwritten{\in}U(k)$.

La première assertion est claire sur la forme 1.5. (i bis).

L'équivalence de (i) à (iv) est une conséquence bien connue du \uline{Main Theorem} de Zariski\handwritten{.} \handwritten{On} not\handwritten{era} que \deleted{\unsure{la}} l'hypothèse que $\handwritten{\psi}$ est birationnel implique que \deleted{\unsure{$G/$ $X$, donc $G/N$, $X$ est,}} $X$ est génériquement réduit, donc (comme $p:X\handwritten{\to}G/N$ est fidèlement plat) que $G/N$ est génériquement réduit donc réduit\handwritten{;} par suite si $W$ est réduit, il en \deleted{\unsure{e}}est de même de $X$ qui est \deleted{fibré sur $G$} plat sur $G/N$ à fibres \unsure{géométriques|géométriqques} isomorphes\deleted{\unsure{à}} à $W$ donc réduites. En tous cas, que \handwritten{$W$} s\handwritten{oi}t réduit ou non, $X$ est irréduct\handwritten{i}ble (comme image de $\deleted{\unsure{X|G}}G\mathrm{x}W$ qui l'est, $G$ et $W$ l'étant), et l'équivalence de (i) à (iii) résulte alors du Main T\deleted{\unsure{he}}heorem appliqué au morphisme $X_{\text{réd}}\handwritten{\to}V$ induit par $\deleted{\phi}\handwritten{\psi}$.

\uline{Remarque} 1.7. Sous les conditions de 1.6. , un point $a$ de $V(k)$ qui satisfait aux conditions équivalentes (i) à (iii) pourra être appelé un \uline{point} \deleted{\unsure{W-rgéu}} $W$-\uline{régulier} de $V$ (relatif aux opérations de $G$ sur $V$). \deleted{Cette notion manifestement est stable par extension du corps de base $k\quad k'$, avec $k'$ algébriquement clos.} Si on ne suppose plus $k$ algébriquement clos, on dira encore qu\deleted{e} 'un point $a$ de $V$ est \inserted{$W$-}régulier si $\phi^{-1}(a)$ est \inserted{non vide et} fini sur $k(a)$\handwritten{.} \deleted{\handwritten{\unsure{i.e.}\ \illegible{3}}} Sous les conditions plus générales de 1.5. (i) à (i ter) , i.e. $\handwritten{\psi}$ étant génériquement étale (mais pas nécessairement birationnel), \deleted{\unsure{on}}il y a lieu d'introduire le degré de projection $d$ de $\handwritten{\psi}$, et le plus grand ouvert $U$ de $V$ tel que $\deleted{\phi}\handwritten{\psi}$ induise un morphisme $\deleted{\phi}\handwritten{\psi}^{-1}(U)\handwritten{\to}U$ qui soit \uline{étale} et \uline{fini}. Lorsque $V$ est normal\inserted{\handwritten{et} $X$ réduit (p.ex. $W$ réduit)}\handwritten{,} alors un $a\handwritten{\in}V$ est dans $U$ si et seulement si $\deleted{\phi}\handwritten{\psi}^{-1}(a)$ est fini et a un degré séparable égal à $d$ , i.e. (lorsque $k$ est algébriquement clos et $a$ est ra\handwritten{t}io\handwritten{nn}el
% archive sheet number: 6
